% Transcription — fonds Alexandre Grothendieck, shelfmark 47, batch 1 (pages 1-10).
% Established from the facsimile archives/batches/47/batch-01.pdf (a mirror of
% https://grothendieck.umontpellier.fr/47.pdf), per the skill
% transcribe-grothendieck. The folder is ten pages and this is its only batch.
%
% Page 1 is a cover sheet carrying, in his hand, the title the inventory
% reproduces — "Produit de convolution sur les VA et transformation de Fourier
% dans les VA" — and nothing else: it takes no \page{} and its words open the
% file as the section heading. Pages 2 to 10 are one continuous run in blue
% ink, paginated 2-10 by the archivists only, and every one of them is
% mathematics; no page is skipped.
%
% What the run is: a Fourier transform for abelian schemes over an arbitrary
% base, stated on the derived category of perfect complexes. Proposition 1 (pp.
% 2-3) proves that the kernel transform Phi_L exchanges Pontrjagin convolution
% and tensor product, by base change, projection formula and the theorem of
% the cube. The theorem of pp. 3-4 states that Phi_{sL} composed with Phi_L is
% (-1)^n s_A^* up to a twist tau_A[-n], hence an equivalence; its proof (pp.
% 4-5) reduces the kernel to a pushforward of the Weil sheaf, computed by Lemma
% 1, which pp. 6-7 deduce from a vanishing lemma he attributes to Mumford and
% admits. Remarks 1-5 (pp. 7-10) give the images of skyscrapers and degree-zero
% line bundles, rank and Euler characteristic, a heuristic formula for the
% Chern classes of the transform of a line bundle with the divisibility it
% would imply, the Chow-theoretic version via Riemann-Roch (with a Todd class
% correction and a sign (-1)^n he says Kleiman's "Dix exposés" formulas
% missed), duality and transport of structure.
%
% Notation fixed for the batch: his blackboard R and L for derived functors
% are \mathbb{R}, \mathbb{L}; the L over * and over \otimes for the derived
% operations is \overset{L}{*}, \overset{L}{\otimes}; the hat on Lambda (top
% exterior power) is \hat\Lambda, once noted; the sheaf on B x A deduced from
% L by swapping factors, which he writes with a superscript s before L, is
% {}^{s}L. A group he writes "Cw" — with a raised dot, and subscript Q — on
% pages 8-10 is transcribed as the glyphs read, \mathrm{Cw}, and flagged
% uncertain at its first occurrence: it is where the Chow ring would stand.
%
% The hand is fast; the formulas carry throughout and most of the prose does,
% but the margins of pages 2 and 4 are written slanting up the left edge and
% are only partly legible, and several passages are reworked with the earlier
% attempts struck. Those are transcribed as \struck{} where they can be read.
%
% Pass: Fable 5.1 (claude-fable-5-1), 2026-09-03 — first pass, unchecked against the pages by a human.
\documentclass[11pt,a4paper]{article}
% Le préambule commun à toutes les transcriptions du fonds Grothendieck.
%
% Il tient en une idée : **l'incertitude se compose, elle ne se résout pas.**
% Une transcription qui lisse un mot douteux a détruit la seule chose pour
% laquelle on la faisait — un lecteur doit pouvoir savoir, à chaque endroit,
% ce qui était sur le feuillet et ce qui a été ajouté. Les six macros qui
% suivent sont donc l'appareil critique, et non de la décoration.
%
% Les mêmes commandes sont reconnues par `scripts/render.mjs`, qui produit la
% vue de lecture du site. Une macro ajoutée ici doit l'être aussi là-bas,
% faute de quoi le rendu échouera bruyamment — ce qui est voulu : un rendu
% silencieusement faux est pire qu'un rendu absent.

\ProvidesPackage{grothendieck}[2026/08/08 Transcriptions du fonds Grothendieck]

\usepackage{amsmath,amssymb,amsthm}
\usepackage{mathrsfs}
\usepackage{stmaryrd}
% Les diagrammes commutatifs sont partout dans ce fonds : c'est la langue
% même de la géométrie algébrique de Grothendieck, et une transcription qui
% les rendrait en prose ne transcrirait rien.
\usepackage{tikz-cd}
% Français par défaut — c'est la langue des feuillets — mais l'anglais est
% chargé aussi : la traduction et le résumé sont des documents anglais, et la
% typographie française (espaces avant les deux-points) y serait une faute.
% Ces documents basculent par \selectlanguage{english} en tête de corps.
\usepackage[english,french]{babel}
\usepackage{fontspec}
\usepackage{geometry}
\geometry{a4paper,margin=25mm,marginparwidth=28mm}
\usepackage{marginnote}
\usepackage{xcolor}
% ulem, not soul. `soul`'s \st and \ul are built on a character-by-character
% scan that XeLaTeX's font machinery breaks: the document fails with an
% undefined control sequence rather than degrading. ulem does the same two jobs
% and is Unicode-engine safe. `normalem` stops it hijacking \emph, which the
% transcriptions use for Grothendieck's own emphasis.
\usepackage[normalem]{ulem}
\usepackage{hyperref}
\hypersetup{colorlinks=true,linkcolor=black,urlcolor=black,pdfborder={0 0 0}}

% --- Métadonnées du lot -----------------------------------------------------
% Reprises telles quelles par le script de rendu, qui les affiche en tête de
% la vue de lecture. Elles fixent ce que le fichier prétend couvrir : sans
% elles, une transcription flotte sans point d'attache dans un fonds de seize
% mille feuillets.

\newcommand{\folder}[1]{\gdef\grothFolder{#1}}
\newcommand{\batch}[1]{\gdef\grothBatch{#1}}
\newcommand{\leaves}[2]{\gdef\grothFirst{#1}\gdef\grothLast{#2}}
\newcommand{\foldertitle}[1]{\gdef\grothFoldertitle{#1}}
\gdef\grothFolder{?}\gdef\grothBatch{?}\gdef\grothFirst{?}\gdef\grothLast{?}\gdef\grothFoldertitle{}

% --- Le résumé --------------------------------------------------------------
%
% Il ouvre la lecture modernisée, sous le titre « Résumé », et il est composé
% en petit. Ce n'est pas un ornement : c'est le seul passage du document qui
% ne soit pas des mathématiques, et un lecteur doit pouvoir le reconnaître
% avant de l'avoir lu — puis le sauter s'il connaît déjà le sujet, ou s'y
% arrêter s'il ne le connaît pas. Le filet qui le suit marque la reprise du
% corps.

\newenvironment{resume}
  {\small\setlength{\parskip}{0.45em}}
  {\par\medskip\hrule\medskip}

% --- Mots-clés ---------------------------------------------------------------
%
% En anglais, en clôture du résumé de la lecture modernisée : le vocabulaire
% moderne sous lequel chercher ce que les feuillets construisent. Le manifeste
% du site (`npm run manifest`) les extrait pour étiqueter la cote entière —
% cette ligne est la source unique des tags, il n'y a pas de fichier à côté.
\newcommand{\keywords}[1]{\par\smallskip\noindent
  {\footnotesize\textsc{Keywords} --- \textit{#1}}\par}

% --- L'appareil critique ----------------------------------------------------

% Le numéro de feuillet, en marge. C'est l'ancre qui permet de régler un
% désaccord sur une formule en regardant *un* feuillet plutôt qu'une chemise
% de six cents. Le site s'en sert aussi pour tourner les pages du fac-similé
% quand on fait défiler la transcription.
\newcommand{\leaf}[1]{%
  \marginnote{\footnotesize\color{gray!70}\textsf{#1}}%
  \label{leaf:#1}\ignorespaces}

% Illisible : on ne devine pas. Un mot inventé qui se lit comme les autres est
% le pire résultat possible.
\newcommand{\ill}{\textcolor{red!60!black}{[\,\ldots\,]}}

% Lecture douteuse : le mot est donné, et signalé comme incertain.
\newcommand{\uncertain}[1]{\uline{#1}}

% Ajout de l'éditeur — une lettre manquante, un mot sous-entendu.
\newcommand{\add}[1]{\textcolor{blue!50!black}{[#1]}}

% Biffé par Grothendieck lui-même. Ce qu'il a rayé fait partie du document :
% une rature dit souvent où la pensée a changé de direction.
\newcommand{\struck}[1]{\sout{#1}}

% Note du transcripteur, sur l'état matériel du feuillet.
\newcommand{\note}[1]{\par\smallskip{\small\color{gray!60!black}\textit{[#1]}}\par\smallskip}

% Note marginale de Grothendieck, distincte du corps du texte.
\newcommand{\marginal}[1]{\par\smallskip\noindent\hspace{1em}%
  {\small\color{gray!60!black}#1}\par\smallskip}

% --- Titre ------------------------------------------------------------------

\AtBeginDocument{%
  \begin{center}
    {\large\bfseries Fonds Alexandre Grothendieck --- cote n\textsuperscript{o}~\grothFolder}\\[2pt]
    {\itshape\grothFoldertitle}\\[4pt]
    {\small feuillets \grothFirst--\grothLast\ (lot \grothBatch)}\\[2pt]
    {\footnotesize\color{gray!60!black}Transcription établie d'après le fac-similé numérisé
     par l'Université de Montpellier.\\
     Les lectures incertaines sont soulignées, les illisibles notées [\,\ldots\,],
     les ajouts entre crochets.}
  \end{center}
  \vspace{1em}}

\endinput


\folder{47}
\batch{1}
\pages{1}{10}
\dating{s.d. — le groupe « Variétés abéliennes » (45 à 56) est daté 1961-1973}
\watermark{Édition de démonstration}
\foldertitle{Produit de convolution sur les VA et transformation de Fourier dans les VA : notes manuscrites (s.d.)}

\begin{document}

\section*{Produit de convolution sur les VA et transformation de Fourier dans les VA}

\note{titre de sa main sur le premier feuillet, qui ne porte rien d'autre et
ne reçoit donc pas de \texttt{page} ; « VA » abrège variétés abéliennes. Les
neuf feuillets qui suivent forment une suite continue}

\page{2}

\textbf{Proposition 1.} Soient $S$ un schéma, $A$ un schéma abélien sur $S$,
$X$ un schéma sur $S$, $L$ un \struck{\ill{}} faisceau inversible sur
\struck{\ill{}} $A \times_S X$, \struck{Considérons l'homomorphisme} trivial
sur $e_A \times_S X$, et tel que \struck{pour tout \ill{}} pour tt $x \in X$,
le faisceau $L_x$ sur $A_x = A_s \otimes_{k(s)} k(x)$ (où $s$ est
\uncertain{l'image} de $x$) soit alg.\ équivalent à zéro. Considérons
l'application
\[
  \varphi_L \colon K^{\cdot}(A) \longrightarrow K^{\cdot}(X)
\]
définie par \struck{l'application additive} le foncteur de catégories
triangulées $\Phi \colon D_{\mathrm{parf}}(A) \to D_{\mathrm{parf}}(X)$
\[
  \Phi_L(F) = \mathbb{R}pr_{2*}\bigl(\mathbb{L}pr_1^{*}(F) \otimes L\bigr)
\]
Alors $\varphi_L$ transforme produit de Pontrjagin $*$ en produit ordinaire.
Plus précisément, on a un isomorphisme de foncteurs de catégories triangulées
\[
  \Phi_L\bigl(F \overset{L}{*} G\bigr) \;\simeq\;
  \Phi_L(F) \overset{L}{\otimes} \Phi_L(G)
  \qquad (F, G \in \mathrm{Ob}\, D_{\mathrm{parf}}(A)).
\]

\marginal{peut se généraliser au cas $A$ un schéma en groupes, $L$ sur
$A \times_S X$ satisfaisant à
$(\pi \times \mathrm{id}_X)^{*}(L) \simeq pr_{13}^{*}(L) \otimes pr_{23}^{*}(L)$,
et $F$ et $G$ sur $A$ \ill{} $\in \mathrm{Ob}\, D^{b}_{\mathrm{qcoh}}(A)$
\ill{} ; on peut aussi généraliser au cas \ill{} concernant des complexes sur
des $C$-torseurs, $P \times Q \to R$, \uncertain{$C \times Q \to C$}, \ill{}
et des faisceaux $L$, $L'$, $L''$ sur $R \times X$, $P \times X$, $Q \times X$
\ill{} tels que
$(\pi \times \mathrm{id}_X)^{*}(L) \simeq pr_{13}^{*}(L')\, pr_{23}^{*}(L'')$}

\note{la note marginale court en biais le long du bord gauche du feuillet,
en deux blocs ; entre les deux, plusieurs lignes ne se laissent pas lire}

Si le th.\ est démontré dans la situation donnée, il en résulte qu'il est vrai
\uncertain{aussi} dans la situation $(A, X', L')$ qui s'en déduit par image
inverse à l'aide d'un $S$-morphisme $f \colon X' \to X$, car
$\Phi_{L'} \simeq \mathbb{L}f^{*}(\Phi_L)$, et $\mathbb{L}f^{*}$ commute à
$\overset{L}{\otimes}$. Soient donc $B$ le schéma abélien dual de $A$, et
$L_0$ le faisceau de Weil sur $A \times_S B$. Alors l'hypothèse sur $L$
implique que $L$ est $\simeq$ l'image inverse de $L_0$ par un morphisme
$X \to B$. Donc on est ramené au cas $X = B$, $L = L_0$. \struck{On} Or on a
des isom.\ fonctoriels
\[
  (*) \qquad F \overset{L}{*} G = \mathbb{R}\pi_{*}\bigl(F \overset{L}{\otimes}_S G\bigr)
  \quad \text{et}
\]

\page{3}

\begin{tikzcd}[column sep=small]
A \times A \arrow[d, "\pi"'] & A \times_S A \times_S B \arrow[l, "pr_{12}"'] \arrow[d, "\pi \times \mathrm{id}_B"] \arrow[dr, "pr_3"] & \\
A & A \times_S B \arrow[l, "pr_1"] \arrow[r, "pr_2"'] & B
\end{tikzcd}

\note{« cart » est inscrit dans le carré de gauche}

\begin{tikzcd}[column sep=small]
A \times_S A \times_S B \arrow[r, "pr_3"] \arrow[d, "f"'] \arrow[dd, bend right=70, "pr_{12}"'] & B \arrow[d, "\Delta_B"] \\
A \times_S A \times_S B \times_S B \arrow[r, "pr_{34}"] \arrow[d, "pr'_{12}"] & B \times_S B \\
A \times_S A &
\end{tikzcd}

\note{« cart » est inscrit dans le carré ; la flèche verticale $f$ porte sur
la page $\mathrm{id}_{A \times A} \times \Delta_B = f$}

\begin{align*}
\Phi\bigl(F \overset{L}{*} G\bigr)
  &\underset{\text{déf}}{\simeq}
    \mathbb{R}pr_{2*}\bigl(\mathbb{L}pr_1^{*}(F \overset{L}{*} G) \otimes L\bigr)
  \underset{(*)}{=}
    \mathbb{R}pr_{2*}\bigl(\mathbb{L}pr_1^{*}(\pi_{*}(F \overset{L}{\otimes}_S G)) \otimes L\bigr) \\
  &\underset{\text{chgt. de base}}{\simeq}
    \mathbb{R}pr_{2*}\bigl(\mathbb{R}(\pi \times \mathrm{id}_B)_{*}(\mathbb{L}pr_{12}^{*}(F \overset{L}{\otimes}_S G)) \otimes L\bigr) \\
  &\underset{\text{formule de proj.}}{\simeq}
    \mathbb{R}pr_{2*}\bigl(\mathbb{R}(\pi \times \mathrm{id}_B)_{*}(\mathbb{L}pr_{12}^{*}(F \overset{L}{\otimes}_S G) \\
  &\hphantom{\underset{\text{formule de proj.}}{\simeq}\mathbb{R}pr_{2*}\bigl(}
    \otimes (\pi \times \mathrm{id}_B)^{*}(L))\bigr) \\
  &\underset{\text{transit.}}{\simeq}
    \mathbb{R}pr_{3*}\bigl(\mathbb{L}pr_{12}^{*}(F \overset{L}{\otimes}_S G) \otimes (\pi \times \mathrm{id}_B)^{*}(L)\bigr)
\end{align*}

\note{le « $\otimes L$ » de la première ligne est ajouté au-dessus, oublié
d'abord}

Or on a (th.\ des carrés) un isom.\ can.
\[
  (\pi \times \mathrm{id}_B)^{*}(L) \simeq pr_{13}^{*}(L)\, pr_{23}^{*}(L)
\]
où $pr_{13}, pr_{23} \colon A \times A \times B \to A \times B$ sont
respectivement $pr_1 \times \mathrm{id}_B$ et $pr_2 \times \mathrm{id}_B$. On
en conclut
\[
  \mathbb{L}pr_{12}^{*}(F \overset{L}{\otimes}_S G) \otimes (\pi \times \mathrm{id}_B)^{*}(L)
  \;\simeq\;
  \mathbb{L}f^{*}\bigl(pr^{\prime *}_{12}(F \overset{L}{\otimes}_S G) \otimes pr^{\prime *}_{13}(L) \otimes pr^{\prime *}_{24}(L)\bigr)
\]
où $pr'_{13}$, $pr'_{24}$ sont les deux projections
$A \times_S A \times_S B \times_S B \to A \times_S B$. On a d'ailleurs
\struck{on a} (\uncertain{ass.\ du $\otimes$})
\[
  pr^{\prime *}_{12}(F \overset{L}{\otimes}_S G) \otimes pr^{\prime *}_{13}(L) \otimes pr^{\prime *}_{24}(L)
  \;\simeq\;
  pr^{\prime *}_{13}(F \otimes L) \overset{L}{\otimes} pr^{\prime *}_{24}(G \otimes L)
\]
et on en conclut
\begin{align*}
\Phi\bigl(F \overset{L}{*} G\bigr)
  &\simeq \mathbb{R}pr_{3*}\bigl(\mathbb{L}f^{*}(\quad)\bigr)
  \underset{\text{chgt. de base}}{\simeq}
    \mathbb{L}\Delta_B^{*}\Bigl(\mathbb{R}pr_{34*}\bigl(pr^{\prime *}_{13}(F \otimes_S L) \overset{L}{\otimes} pr^{\prime *}_{24}(G \otimes_S L)\bigr)\Bigr) \\
  &\underset{\text{Künneth}}{\simeq}
    \mathbb{L}\Delta_B^{*}\bigl(\mathbb{R}pr_{2*}(F \otimes L) \overset{L}{\otimes}_S \mathbb{R}pr_{2*}(G \otimes L)\bigr) \\
  &\underset{\text{déf}}{\simeq}
    \mathbb{L}\Delta_B^{*}\bigl(\Phi(F) \overset{L}{\otimes}_S \Phi(G)\bigr)
  \simeq \Phi(F) \overset{L}{\otimes} \Phi(G)
\end{align*}
cqfd.

\note{les parenthèses vides de $\mathbb{L}f^{*}(\quad)$ sont sur la page : un
trait les relie à l'expression soulignée deux lignes plus haut}

\textbf{Théorème \struck{Proposition 2}.} Soient $A$, $B$ deux schémas
abéliens sur $S$, de dim relative $n$, $L$ un Module inversible birigidifié
sur $A \times B$, définissant une dualité entre $A$ et $B$.
\struck{Considérons} Alors le composé \struck{des trans\ill{}}
\[
  K(A) \xrightarrow{\;\varphi_L\;} K(B) \xrightarrow{\;\varphi_{{}^{s}L}\;} K(A)
\]
est $(-1)^n s_A^{*}$ (où $s_A$ est la symétrie de $A$) (donc, par exemple,
$\varphi_L$ et $(-1)^n s_A^{*} \varphi_{{}^{s}L}$ sont des isomorphismes
inverses l'un de l'autre).

\note{« de dim relative $n$ » est ajouté en interligne ; ${}^{s}L$ est le
faisceau sur $B \times A$ déduit de $L$ par la symétrie des facteurs ; le
$(-1)^n$ de la dernière parenthèse est ajouté au-dessus de la ligne. La suite
cite cet énoncé comme « prop.\ 2 »}

\page{4}

Plus précisément, le foncteur composé
\[
  D_{\mathrm{parf}}(A) \xrightarrow{\;\Phi_L\;} D_{\mathrm{parf}}(B)
  \xrightarrow{\;\Phi_{{}^{s}L}\;} D_{\mathrm{parf}}(A)
\]
est \struck{isomorphe au foncteur} can.\ isomorphe au foncteur
$\mathbb{L}s_A^{*} \simeq \mathbb{R}s_{A*}$ suivi de $\tau_A[-n]$,
\struck{\ill{}} où $\tau_A \simeq \hat\Lambda\, t_A \simeq \hat\Lambda\, t_B$
\struck{respectivement}. Donc $\Phi_L$ et $\mathbb{L}s_A^{*}\Phi_{{}^{s}L}$
($= \mathbb{R}s_{A*}\Phi_{{}^{s}L}$) \struck{\ill{}} sont des équivalences
des catégories triangulées de $D_{\mathrm{parf}}(A)$ et $D_{\mathrm{parf}}(B)$.

\note{$\hat\Lambda$ : le $\Lambda$ porte un accent, remplacé plus bas
(p.\ 6) par un $n$ ; c'est la puissance extérieure maximale. $t_A$ est le
faisceau tangent de $A$ le long de la section nulle, cf.\ p.\ 6}

\marginal{inutile de se borner aux complexes parfaits, prendre \ill{}
p.\ ex.\ dans $D^{b}_{\mathrm{qcoh}}$ \struck{\ill{}}}

\struck{\ill{}}

\note{un mot encerclé et biffé, en haut à droite, ne se lit pas}

\textbf{Corollaire \uncertain{2}.} \struck{\ill{}} $\Phi_L$ transforme
produit ordinaire en produit de convolution :
\[
  \Phi_L\bigl(F \overset{L}{\otimes} G\bigr) \otimes \omega_A[n]
  \;\simeq\;
  \bigl[\Phi_L(F) \otimes \omega_A[n]\bigr] * \bigl[\Phi_L(G) \otimes \omega_A[n]\bigr]
\]
i.e.
\[
  \Phi_L\bigl(F \overset{L}{\otimes} G\bigr)
  \;\simeq\;
  \bigl[\Phi_L(F) * \Phi_L(G)\bigr] \otimes \omega_A[n].
\]

\note{le début de l'énoncé est remanié et largement biffé ; on lit encore,
biffés, « $\Phi_L^{\omega_A[n]}$ », « $\overset{L}{\otimes}$ en
$\overset{L}{*}$ » et « $x \mapsto (-1)^n\, \mathrm{cl}\,\omega_A\,
\varphi_L(x)$ », et, non biffé, en interligne au-dessus, « $F \mapsto
\Phi_L(F \otimes \omega_A[n]) \simeq \Phi_L(F) \otimes \omega_A[n]$
transforme »}

Prouvons la prop.\ 2. Il suffit de définir un isom.\
$\Phi_{{}^{s}L} \cdot \Phi_L = \mathrm{id}_{D_{\mathrm{parf}}(A)}$. Or le
composé des foncteurs est de la forme
$F \mapsto \mathbb{R}pr_{2*}\bigl(pr_1^{*}(F) \overset{L}{\otimes} M\bigr)$,
où $M \in D_{\mathrm{parf}}(A)$ est donné par
\[
  M = \mathbb{R}pr_{13*}\bigl(pr_{12}^{*}(L)\, pr_{23}^{*}({}^{s}L)\bigr)
\]

\marginal{$A \times B \times A$}

\note{« $D_{\mathrm{parf}}(A)$ » est ce qui est écrit ; le noyau $M$ vit sur
$A \times A$. La note marginale nomme l'espace des projections}

Il faut donc calculer ce $\mathbb{R}pr_{13*}$, et pour ceci, posons
$C = A \times A$, et considérons l'isom.\ $A \times B \times A \simeq$
\struck{\ill{}} $C \times B$. Le module inversible
$N = pr_{12}^{*}(L)\, pr_{23}^{*}({}^{s}L)$ sur $C \times B$ est
\struck{\ill{}} birigidifié de façon canonique, et correspond à un
homomorphisme $C \xrightarrow{\;\Psi\;} B^{*}$, qui n'est autre que l'hom.\
de $A \times A$ dans $B^{*}$ dont les deux composantes sont égales à
$\varphi \colon A \to B$ (l'hom.\ de $A$ dans $B^{*}$ défini par $L$).

\note{« deux » est ajouté en interligne}

\page{5}

\begin{tikzcd}
C \times_S B \arrow[r, "\Psi \times \mathrm{id}_B"] \arrow[d, "pr_1"'] & B^{*} \times_S B \arrow[d, "pr'_1"] \\
C \arrow[r, "\Psi"] & B^{*}
\end{tikzcd}

\note{« cart » dans le carré ; le second facteur du coin supérieur droit
porte un astérisque biffé}

Ceci dit, on aura \struck{$\mathbb{R}pr_{2*}(N)$}
$N \simeq (\Psi \times \mathrm{id}_B)^{*}(W_B)$ ($W_B$ le Module de Weil sur
$B \times_S B^{*}$), d'où
\[
  \mathbb{R}pr_{2*}(N) \simeq \mathbb{R}pr_{2*}(\Psi \times \mathrm{id}_B)^{*}(W_B)
  \underset{\text{chgt de base}}{\simeq}
  \mathbb{L}\Psi^{*}\bigl(\mathbb{R}pr'_{2*}(W_B)\bigr).
\]

\note{la ligne écrit $pr_2$, $pr'_2$ là où le diagramme et le corollaire
ci-dessous ont $pr_1$, $pr'_1$}

\textbf{Lemme 1.} Soient $B$ un schéma abélien sur $S$, de dim.\ relative $n$,
$W$ le faisceau de Weil sur $B^{*} \times_S B$,
$p \colon B^{*} \times_S B \to B^{*}$ la projection. Alors
\[
  \mathbb{R}p_{*}(W) \simeq e_{*}\bigl(\struck{\ill{}}\; \tau_{B^{*}}[-n]\bigr)
\]
($\tau_{B^{*}} = \hat\Lambda^{n}\, t_{B^{*}}$, $t_{B^{*}}$ le faisceau
tangent à $B^{*}$ le long de la section $0$) (isom.\ canonique).

\note{« de dim.\ relative $n$ » est ajouté en interligne, relié par une
accolade ; les astérisques des deux $B$ de l'énoncé sont retouchés, un ajouté,
un biffé, et la lecture $B^{*} \times_S B$ suit le diagramme}

\marginal{$\simeq e_{*}(\omega^{-1})$ où
$\omega = \uncertain{\det{}^{*}}(\Omega^{1}_{B^{*}/S})$}

On en conclut

\textbf{Corollaire.} Soient $C$, $B$ schémas abéliens sur $S$, $N$ un Module
inversible birigidifié sur $C \times_S B$, $\Psi \colon C \to B^{*}$ l'hom.\
correspondant. Alors on a
\[
  \mathbb{R}pr_{1*}(N) \simeq \mathbb{L}\Psi^{*}\bigl(\uncertain{e_{*}\tau_{B^{*}}}[-n]\bigr)
\]

\note{les symboles entre la parenthèse et $[-n]$ sont surchargés et ne se
lisent qu'avec doute}

En particulier, si $\Psi$ est un épimorphisme, \struck{donc} plat, et si
$K = \operatorname{Ker}\Psi$, donc $K$ est un sous-schéma en groupes de $C$
plat et de prés.\ finie sur $S$ donc \struck{défini} \struck{$\mathcal{O}_K$
\ill{}} \uncertain{parfait} sur $C$, et on a
\[
  \mathbb{R}pr_{1*}(N) \simeq \mathcal{O}_K[-n] \otimes_S \tau_{B^{*}}
\]

Dans le cas qui nous occupe, $C = A \times A$, \struck{\ill{}},
$\varphi \colon A \simeq \uncertain{C^{*}}$, et \struck{\ill{}} par suite $K$
est l'antidiagonale de $A \times_S A$, noyau de
$\pi_A \colon A \times_S A \to A$, qui est le graphe de la symétrie $s_A$ de
$A$. Donc \struck{\ill{}} \uncertain{modulo dém.\ du} Lemme 1.

\note{la lettre lue $C^{*}$ est nette ; l'argument demande $B^{*}$}

\page{6}

Démonstration du Lemme 1. Il sera une conséquence du

\textbf{Lemme 2 (Mumford).} Soit $B$ schéma abélien sur un corps (alg.\ clos)
$k$, et $L$ un Module inv.\ sur $B$ alg.\ équivalent à zéro, \struck{Alors}
et non isomorphe à $\mathcal{O}_B$. Alors on a
\[
  H^{i}(X, L) = 0 \quad \text{pour tout } i
\]

\note{« $X$ » est ce qui est écrit, pour $B$. Suit un encadré entièrement
biffé, où l'on distingue $H^{0}(X, L) = 0$, $k$, $L \simeq \mathcal{O}$}

\struck{Démonstration du Lemme 2. L'assertion sur les $H^{0}$ est triviale,
et il reste à prouver que $H^{i}(X, L) = 0$ \ill{}}

Nous allons ici admettre le lemme 2, et nous montrons comment on
\struck{conclut} déduit le lemme 1. Nous voyons déjà grâce \uncertain{au}
lemme 2 que $\mathbb{R}p_{*}(W)$ est concentré sur la section unité de
$B^{*}$. Nous allons définir un homomorphisme canonique (on écrit
$e_{B^{*}} = e$)
\[
  \mathbb{R}p_{*}(W) \longrightarrow \struck{\ill{}}\; e_{*}(\mathcal{O}_S) \otimes_{\mathcal{O}_S} \tau_{B^{*}}[-n]
\]
(comme \ill{}, \ill{} $R^{i}p_{*}(W) = 0$ si $i > n$) ou, ce qui revient au
même, un homomorphisme canonique
\[
  R^{n}p_{*}(W) \longrightarrow \struck{\otimes_S}\; \tau_{B^{*}}
\]
ou encore un hom.
\[
  \struck{e_{B}}\; e^{*}\bigl(R^{n}p_{*}(W)\bigr) \longrightarrow \struck{\otimes_S}\; \tau_{B^{*}}
\]

\begin{tikzcd}
B^{*} \times_S B \arrow[d, "p"'] & B \arrow[l, "e'"'] \arrow[d, "g"] \\
B^{*} & S \arrow[l, "e"]
\end{tikzcd}

Or on a \ill{} (formule de changement de base, comme
$e^{\prime *}(W) \simeq \mathcal{O}_B$, et \struck{\ill{}}) définissant un
hom.
\[
  R^{n}g_{*}(\mathcal{O}_B) \longrightarrow \struck{\otimes_S}\; \tau_{B^{*}}
\]
Or on sait que l'on a un isom canonique
\[
  R^{n}g_{*}(\mathcal{O}_B) \simeq \hat\Lambda\, R^{1}g_{*}(\mathcal{O}_B) \simeq \hat\Lambda^{n}\, t_{B^{*}}
\]
où $t_{B^{*}}$ est le faisceau tangent de $B^{*}$ le long de $e$. On prend
alors \struck{\ill{}} l'hom identique $\hat\Lambda^{n}\, t_{B^{*}} \simeq \tau_{B^{*}}$.

\page{7}

Ceci posé, je dis que l'homomorphisme précédent
\[
  \mathbb{R}p_{*}(W) \longrightarrow e_{*}(\mathcal{O}_S) \otimes \tau_{B^{*}}[-n]
\]
est un isomorphisme. Il suffit pour ceci de prouver que l'homomorphisme
induit \struck{pour les} en tout point \struck{$x$} $b \in B^{*}$ par
$\mathbb{L}i_b^{*}$ est un isomorphisme (histoire de \ill{}-cylindres
\ldots). En les points $b \in B^{*} - e(S)$, ce n'est autre que le lemme 2.
En un point \struck{$b \in e$} \uncertain{$b$ de} $e(S)$, on est ramené au
cas où $S$ est le spectre d'un corps algébriquement clos, à faire une
vérification de compatibilité \ldots\ldots

\textbf{Remarques} 1) a) $\Phi_L$ transforme $e_A$ en $1_B$, $e_s$ ($s$
section de $A/S$) en $L_s$ (faisceau inv.\ alg.\ $\simeq 0$ sur $B$ rel
$/S$), donc $e_s * F$ en $L_s \otimes \Phi_L(F)$.

b) $\Phi_L$ transforme \struck{$\omega_A[n]$ en \ill{}, $\omega_A[n]$} $L_t$
($t$ section de \uncertain{$B$} sur $S$) en $\tau[-n]\, e_{-t}$, i.e.\
$L_t\, \omega[n]$ en $e_{-t}$, \struck{donc} $F \otimes L_t$ \struck{\ill{}}
en $\bigl[\Phi_L(F) * e_{-t}\bigr]$ \struck{\ill{}}.

\note{la lettre lue $B$ est formée comme un $D$, ici comme p.\ 5 dans
$W_B$ ; $e_A$, $e_s$, $e_{-t}$ sont les faisceaux structuraux des sections
correspondantes}

c) $\varepsilon_B(\Phi_L(F)) = \chi_A(F)$,
$\chi_B(\Phi_L(F)) = (-1)^n\, \varepsilon_A(F)$.

\marginal{$\varepsilon = $ \ill{} ; égalité dans $H^{0}(S, \mathbf{Z})$}

2) $\det{}^{*} \circ \Phi_L$ définit un homomorphisme
\[
  \varphi \colon \underline{\mathrm{Pic}}_{A/S} \longrightarrow \underline{\mathrm{Pic}}_{B/S}
\]
qui, pour toute section \struck{\ill{}} $\delta$ de $\underline{NS}_{A/S}$,
définit un \ill{}
\[
  \underline{\mathrm{Pic}}^{\delta}_{A/S} \longrightarrow \underline{\mathrm{Pic}}^{-\delta'}_{B/S}
\]
compatible avec \uncertain{$-(-\tilde\delta{}') = \tilde\delta{}'$}

\page{8}

si $\delta' \in \underline{NS}_{B/S}(S)$ été défini comme \ill{}
constitué par ailleurs, par
\[
  \tilde\delta{}'\, \tilde\delta = -\chi(\delta)\, \mathrm{id}_A
\]
\ldots\ldots

On peut se proposer, pour une section $\lambda$ de $\underline{\mathrm{Pic}}_{A/S}$,
définissant $L$ faisceau inv.\ sur $A$, de « calculer » la classe dans $K(B)$
de $\varphi_L(\struck{\ill{}}\, \mathrm{cl}(L))$ en termes des invariants de
la section $\varphi(\lambda)$ de $\underline{\mathrm{Pic}}_{B/S}$, définissant
un faisceau inversible $L'^{-1}$ sur $B$ (de classe $-\delta'$ dans
$NS(B/S)$ si $\delta$ \uncertain{est la classe de} $\lambda$ dans $NS(A/S)$).
Heuristiquement, \struck{\ill{}} devrait avoir (pour $\chi(\delta) \neq 0$
partout sur $S$)
\[
  \varphi_L(\mathrm{cl}(L)) = \chi\; \struck{\ill{}}\; L'^{-1/\chi}
  \qquad \text{où } \chi = \chi(\delta) \in H^{0}(S, \mathbf{Z}).
\]

\marginal{indépendant des résultats de Mumford}

\note{« devrait avoir » est ajouté au-dessus de la ligne, relié à la place
d'un mot biffé ; l'exposant $-1/\chi$ est entre guillemets sur la page}

Au niveau du $\mathrm{ch}_B$, ceci s'écrit
\[
  \mathrm{ch}\, \varphi_L(\uncertain{\mathrm{cl}(L)}) = \chi \exp\Bigl(-\frac{D'}{\chi}\Bigr)
  = \chi - D' + \frac{1}{2\chi} D'^{2} - \frac{1}{3!\,\chi^{2}} D'^{3} \ldots
  + (-1)^n \frac{1}{n!\,\chi^{n-1}} D'^{n},
\]
et au niveau des classes de Chern \struck{\ill{}}
\begin{align*}
  c(\varphi_L(\uncertain{\mathrm{cl}(L)})) &= \Bigl(1 - \frac{D'}{\chi}\Bigr)^{\chi}
  = 1 - D' + \frac{\chi(\chi-1)}{2\chi} D'^{2} - \frac{\chi(\chi-1)(\chi-2)}{3!\,\chi^{2}} D'^{3} + \cdots \\
  &\qquad + (-1)^n \frac{\chi(\chi-1)\cdots(\chi-n+1)}{n!\,\chi^{n-1}} D'^{n}.
\end{align*}

\marginal{$D' \in \uncertain{\mathrm{Cw}^{1}}$ défini par $L'$}

\note{« Cw » est la lecture des glyphes, ici et pp.\ 9-10 : un $C$, une lettre
basse à deux jambages, un point en exposant et l'indice $\mathbf{Q}$ ; c'est
la place de l'anneau de Chow, que la p.\ 10 oppose à une « théorie
cohomologique ». L'argument de $\varphi_L$ dans les deux formules est
surchargé}

? Ces formules sont-elles bien vraies mod torsion ? Elles sont vraies au moins
sur un corps alg.\ clos, mod équivalence homologique (\uncertain{donc le}
$\neq$ \uncertain{car.\ $h$}). Ceci donnerait des propriétés de divisibilité
intéressantes de $\delta'$ : à savoir
\[
  \struck{\ill{}}\; \frac{\chi(\chi-1)\cdots(\chi-k+1)}{k!\,\chi^{k-2}}\, \delta'^{k}
\]
est « entier », i.e.\ provient d'un cycle à coeff.\ entiers.

\note{le point d'interrogation ouvrant le paragraphe est le sien, en marge ;
« équivalence homologique » est souligné deux fois ; la dernière parenthèse
ne se lit qu'avec doute}

\page{9}

3) On vérifie (par R.R.\ pour $A \times A \xrightarrow{\;\pi_A\;} A$) que
l'on a sur $A$
\[
  (*) \qquad \mathrm{ch}(x * y) = \mathrm{ch}(x) * \mathrm{ch}(y)\; \mathrm{Todd}(t_A)
\]
(donc, \ill{} : $\mathrm{ch}(x * y) = \mathrm{ch}(x) * \mathrm{ch}(y)$ si
$\mathrm{Todd}\, t_A = 1$, p.\ ex.\ si $t_A$ est trivial, p.\ ex.\ $S$
local).

Considérons d'autre part l'élément $D$ de $\mathrm{Cw}^{1}(A \times A)$
défini par $L$, et considérons \struck{le} $\exp D$, d'où
\[
  \Psi_D = \varphi_{\exp D} \colon \mathrm{Cw}^{\cdot}(A)_{\mathbf{Q}} \longrightarrow \mathrm{Cw}^{\cdot}(B)_{\mathbf{Q}}.
\]

\note{« $\exp D$ » est écrit au-dessus de « l'élément », relié par un trait}

\struck{Alors \ill{} On a donc, par prop.\ 1 : $\Psi_D(xy) = \Psi_D(x)\Psi_D(y)$.
On veut utiliser prop.\ 2 pour calculer $\Psi_D(a * b)$. On a, si
$a = \mathrm{ch}\, x$, $b = \mathrm{ch}\, y$ : $x * y = \mathrm{ch}\, a *
\mathrm{ch}\, b = \mathrm{ch}(a * b)\, T^{-1}$ ($T = \mathrm{Todd}(t_A)$),
d'où $\Psi_D(x * y) = \Psi_D(\mathrm{ch}(a * b)\, T^{-1})$}

\note{tout le bloc est encadré d'un trait à gauche et barré d'une diagonale ;
c'est une première rédaction, reprise ci-dessous}

\struck{Si on a} \ill{}, par RR pour $A \times_S B \to B$ :
\[
  \mathrm{ch}(\varphi_L(x)) = \Psi_D(\mathrm{ch}(x))\, T \qquad (T = \mathrm{Todd}\, t_A)
\]
\struck{Supposons \ill{}} Posons $a = \mathrm{ch}\, x$, $b = \mathrm{ch}\, y$,
et transformons \ill{}
\[
  \varphi_L(x * y) = \varphi_L(x)\, \varphi_L(y) \qquad \text{(prop 1)}
\]
par $\mathrm{ch}$, on trouve
\[
  \Psi_D(\mathrm{ch}(x * y))\, T = \Psi_D(\mathrm{ch}\, x)\, T\; \Psi_D(\mathrm{ch}\, y)\, T
\]
d'où, \struck{\ill{}} grâce à $(*)$ :
\[
  \Psi_D(\struck{\ill{}}\; \mathrm{ch}\, x * \mathrm{ch}\, y)\, T\, \struck{T^{2}}
  = \Psi_D(\mathrm{ch}\, x)\, \Psi_D(\mathrm{ch}\, y)\, T^{2}
\]
i.e.
\[
  \Psi_D(a * b) = \Psi_D(a)\, \Psi_D(b)
\]

\note{la dernière formule est encadrée ; « $(*)$ » est écrit au-dessus du
facteur biffé de la précédente}

\page{10}

De même, on a (prop.\ 2, cor.)
\[
  \varphi_L(xy) = \bigl(\varphi_L(x) * \varphi_L(y)\bigr)\; \struck{\ill{}}\; (-1)^n \omega_A
\]
d'où en appliquant $\mathrm{ch}$ et $(*)$
\[
  \Psi_D(\mathrm{ch}(xy)) = \struck{\ill{}}\; \Psi_D(\mathrm{ch}\, x) * \Psi_D(\mathrm{ch}\, y)\bigr)\, T\, (-1)^n\, \struck{\exp}\; \mathrm{ch}\, \omega_A
\]
donc comme $\mathrm{ch}(xy) = \mathrm{ch}\, x\; \mathrm{ch}\, y$
\[
  \Psi_D(ab) = \Psi_D(a) * \Psi_D(b)\, (-1)^n\, T\, \mathrm{ch}\, \omega_A
  \qquad
  \begin{cases}
    T = \mathrm{Todd}(t_A) \\
    \omega_A = \hat\Lambda\, t_A^{\vee}
  \end{cases}
\]

\note{la formule est encadrée ; la parenthèse ouvrante manque à la ligne
précédente, comme sur la page}

\marginal{$a, b \in \mathrm{Cw}^{\cdot}(A)_{\mathbf{Q}}$ ; les formules \ill{}
dès que \ill{}}

Ce sont les formules (\struck{\ill{}} \ill{} \uncertain{affirmations}) par
Kleiman dans « 10 Exposés \ldots », sur un corps alg.\ clos et en travaillant
avec une théorie \uncertain{cohomologique} plutôt qu'avec
$\mathrm{Cw}_{\mathbf{Q}}$, --- à cela près que Kleiman n'a pas tenu compte
des facteurs $(-1)^n$ !

4) Il résulte de (\uncertain{th.\ de dualité pour} $A \times_S B \to B$) qu'on
a un \uncertain{isom.}\ fonctoriel
\[
  \Phi_L\bigl(\check F \otimes \omega_A[n]\bigr) \simeq \Phi_L(F)^{\vee}
\]
i.e.
\[
  \Phi_L(F)^{\vee} \simeq \Phi_L\bigl(\check F \otimes \omega_A[n]\bigr)
\]
\uncertain{Donc pour les} \ill{} virtuels, ceci s'écrit
\[
  \varphi_L(x)^{\vee} = (-1)^n\, \omega_A\, \varphi_L(\check x)
\]

\note{la seconde formule est encadrée}

5) Bien sûr, si $u$ est un automorphisme de $A$ et $u'$ l'automorphisme
contragrédient de $B$, on a par transport de structure un isom.\ canonique
\[
  \Phi_L(u.F) \simeq u'.\Phi_L(F)
\]
En particulier, $\Phi_L$ commute à la symétrie $s_A$. Kif-kif pour $\varphi_L$.

\end{document}
